%%%%%%%%%%%%%%%%
%%% exod 24 %%%
%%%%%%%%%%%%%%%%

\Chapter{24}

\Verse{1}
{
	\hE{וְאֶל מֹשֶׁה אָמַר עֲלֵה אֶל יְהֹוָה אַתָּה וְאַהֲרֹן נָדָב וַאֲבִיהוּא וְשִׁבְעִים מִזִּקְנֵי יִשְׂרָאֵל וְהִשְׁתַּחֲוִיתֶם מֵרָחֹק׃}
}
{
	\eE{
		And He said to Moses, “Come up to YHWH: you and Aaron,
		Nadab and Abihu, and seventy of Israel’s elders, and bow
		from a distance.
	}
}
{
%	\footnote{
%		from a distance. This account has a multitude of parallels to the
%		account of the binding of Isaac (Genesis 22): Moses says the same words
%		to the elders that Abraham says to the servant boys: “Sit here … we’ll
%		come back to you.” The Exodus account has servant boys (n‘rîm) as well.
%		Both accounts use the term “from a distance” (mrq). Both use the term
%		“to bow” (hištawôt). Both Moses and Abraham come up a mountain. Both
%		have a burnt offering (ha‘lôt ‘ôlh). The two accounts share a chain of
%		ten verbs: “and he said,” “and he took … and he set,” “and he got up
%		early,” “and he built an altar,” “and he put out his hand,” “and he/it
%		was,” “and he/they got up,” “and he/they came,” “and he/they saw.” All
%		these parallels alert us to the fact that some important connection
%		exists here. The culminating parallels indicate what that connection is:
%		Abraham is rewarded at the end “because you did this thing”; and the
%		people in Exodus promise that “We’ll do all the things.” Abraham is
%		rewarded because “you listened to my voice”; and Exodus reports that
%		“they said with one voice,” and they say, “we’ll listen.” So (1) this is
%		a reminder that the merit of Abraham remains the basis on which all of
%		this now happens. (2) This is a fulfillment of promises made following
%		Abraham’s acts of extraordinary obedience. And (3) Abraham’s obedience
%		is an example for the people to emulate. The contrast is striking as
%		well. Abraham’s obedience is the centerpiece of the first story. In this
%		second story, the people will soon commit their worst act of
%		disobedience, the golden calf event. Note also: The first story is at
%		“the mountain of YHWH.” The second is at “the mountain of God.” And both
%		state that they are about divine appearances. The divine appearances
%		mark the sacredness of the two mountain locations. The “mountain of God”
%		is Horeb/Sinai. The identity of the “mountain of YHWH” is uncertain,
%		although elsewhere it refers to the location of the Temple in Jerusalem.
%		The parallels of these two stories connect the mountain where the
%		greatest revelation takes place to the mountain where the divine promise
%		was confirmed in the past—and where the Temple will stand in the future.
%		(The location of the Temple is identified with the place of the binding
%		of Isaac—Moriah—in 2 Chr Footnote 3:1.) The two sacred mountains are
%		connected; the patriarch and his descendants are connected; past,
%		present, and future are connected—all in the embroidery of the Torah’s
%		wording.
%	}
}

\Verse{2}
{
	\hE{וְנִגַּשׁ מֹשֶׁה לְבַדּוֹ אֶל יְהֹוָה וְהֵם לֹא יִגָּשׁוּ וְהָעָם לֹא יַעֲלוּ עִמּוֹ׃}
}
{
	\eE{
		And Moses will come over alone to YHWH, and they shall not
		come over, and the people shall not come up with him.”
	}
}
{
}

\Verse{3}
{
	\hE{וַיָּבֹא מֹשֶׁה וַיְסַפֵּר לָעָם אֵת כּל דִּבְרֵי יְהֹוָה וְאֵת כּל הַמִּשְׁפָּטִים וַיַּעַן כּל הָעָם קוֹל אֶחָד וַיֹּאמְרוּ כּל הַדְּבָרִים אֲשֶׁר דִּבֶּר יְהֹוָה נַעֲשֶׂה׃}
}
{
	\eE{
		And Moses came and told the people all of YHWH’s words and
		all the judgments. And all the people answered, one voice,
		and they said, “We’ll do all the things that YHWH has
		spoken.”
	}
}
{
}

\Verse{4}
{
	\hE{וַיִּכְתֹּב מֹשֶׁה אֵת כּל דִּבְרֵי יְהֹוָה וַיַּשְׁכֵּם בַּבֹּקֶר וַיִּבֶן מִזְבֵּחַ תַּחַת הָהָר וּשְׁתֵּים עֶשְׂרֵה מַצֵּבָה לִשְׁנֵים עָשָׂר שִׁבְטֵי יִשְׂרָאֵל׃}
}
{
	\eE{
		And Moses wrote all of YHWH’s words. And he got up early
		in the morning and built an altar below the mountain and
		twelve pillars for twelve tribes of Israel.
	}
}
{
}

\Verse{5}
{
	\hE{וַיִּשְׁלַח אֶת נַעֲרֵי בְּנֵי יִשְׂרָאֵל וַיַּעֲלוּ עֹלֹת וַיִּזְבְּחוּ זְבָחִים שְׁלָמִים לַיהֹוָה פָּרִים׃}
}
{
	\eE{
		And he sent young men of the children of Israel, and they
		made burnt offerings, and they made peace-offering
		sacrifices to YHWH: bulls.
	}
}
{
}

\Verse{6}
{
	\hE{וַיִּקַּח מֹשֶׁה חֲצִי הַדָּם וַיָּשֶׂם בָּאַגָּנֹת וַחֲצִי הַדָּם זָרַק עַל הַמִּזְבֵּחַ׃}
}
{
	\eE{
		And Moses took half of the blood and set it in basins and
		threw half of the blood on the altar.
	}
}
{
}

\Verse{7}
{
	\hE{וַיִּקַּח סֵפֶר הַבְּרִית וַיִּקְרָא בְּאזְנֵי הָעָם וַיֹּאמְרוּ כֹּל אֲשֶׁר דִּבֶּר יְהֹוָה נַעֲשֶׂה וְנִשְׁמָע׃}
}
{
	\eE{
		And he took the scroll of the covenant and read in the
		people’s ears. And they said, “We’ll do everything that
		YHWH has spoken, and we’ll listen.”
	}
}
{
%	\footnote{
%		We’ll do … and we’ll listen. It has been noted that the word order here
%		seems backward: doesn’t one have to listen to the command before one can
%		do it?! The promises to do and to listen come together in the order of
%		this sentence in the Hebrew, so it may be that they were a known
%		formula, constituting an oath of absolute obedience. But I understand
%		them to be separate: The people’s promise to do the commandments is a
%		repetition of what they had said before in v. 3. They repeat it formally
%		now because Moses has just read them the written document. The people’s
%		promise to listen is, as Rashbam wrote, a further promise, a commitment
%		to follow the things that their God will say to them in the future. In
%		Hebrew as in English, words for “listening” also imply obeying, as in:
%		“Listen to your mother!” “Do you hear me?!” (As an alternative and more
%		banal explanation of the seemingly backward order: the Septuagint has
%		the two words in reverse order.)
%	}
}

\Verse{8}
{
	\hE{וַיִּקַּח מֹשֶׁה אֶת הַדָּם וַיִּזְרֹק עַל הָעָם וַיֹּאמֶר הִנֵּה דַם הַבְּרִית אֲשֶׁר כָּרַת יְהֹוָה עִמָּכֶם עַל כּל הַדְּבָרִים הָאֵלֶּה׃}
}
{
	\eE{
		And Moses took the blood and threw it on the people, and
		he said, “Here is the blood of the covenant that YHWH has
		made with you regarding all these things.”
	}
}
{
}

\Verse{9}
{
	\hE{וַיַּעַל מֹשֶׁה וְאַהֲרֹן נָדָב וַאֲבִיהוּא וְשִׁבְעִים מִזִּקְנֵי יִשְׂרָאֵל׃}
}
{
	\eE{
		And Moses and Aaron, Nadab and Abihu, and seventy of
		Israel’s elders went up.
	}
}
{
}

\Verse{10}
{
	\hE{וַיִּרְאוּ אֵת אֱלֹהֵי יִשְׂרָאֵל וְתַחַת רַגְלָיו כְּמַעֲשֵׂה לִבְנַת הַסַּפִּיר וּכְעֶצֶם הַשָּׁמַיִם לָטֹהַר׃}
}
{
	\eE{
		And they saw the {\God} of Israel. And below His feet it was like a
		structure of sapphire brick and like the essence of the skies for
		clarity.
	}
}
{
%	\footnote{
%		24:10–11. And they saw {\God}. … And they envisioned {\God}. Some
%		have taken this to mean that they literally, physically see {\God}.
%		But, as Ibn Ezra properly observed, the verb for “they saw” (Hebrew
%		r’h) is also used to refer to seeing in a vision, as when Isaiah
%		says, “I saw the Lord sitting on a throne” (Isa 6:1). This meaning
%		of “they saw” here is confirmed by the parallel use in v. 11 of the
%		word “envisioned” (Hebrew zh), which regularly means to have a
%		vision (as, for example in Isa 1:1). Also, this is consistent with
%		the idea that no one ever sees the actual form of {\God} except
%		Moses on a single occasion. This idea is understood throughout the
%		Tanak, without exception. (On the distinction between vision and
%		regular experience, see The Hidden Face of {\God}, pp. 17–18,
%		62–63.)
%	}
}

\Verse{11}
{
	\hE{וְאֶל אֲצִילֵי בְּנֵי יִשְׂרָאֵל לֹא שָׁלַח יָדוֹ וַיֶּחֱזוּ אֶת הָאֱלֹהִים וַיֹּאכְלוּ וַיִּשְׁתּוּ׃}
}
{
	\eE{
		And He did not put out His hand to the chiefs of the
		children of Israel. And they envisioned {\God}. And they ate
		and drank.
	}
}
{
%	\footnote{
%		He did not put out His hand. Those who take the passage to mean that
%		this group actually sees God understand “He did not put out His hand” to
%		mean that God does not kill them for seeing God. But the text says “He
%		did not put out His hand to the chiefs …” The phrase “to put out one’s
%		hand” followed by the preposition “to” (Hebrew ’el) never refers to
%		killing in the Tanak. (When referring to killing, the phrase occurs with
%		the preposition b.) And, in any case, there are other cases of the
%		expression of God putting out a hand to humans that are meant
%		positively. In Ps 144:7, this refers to God saving someone. And there is
%		an even closer parallel in Jer 1:9, where it takes place in a vision and
%		refers to God reaching out to touch Jeremiah’s mouth to enable him to
%		speak. Here in Exodus, perhaps it means that in their vision they see
%		but are not touched by God.
%	}
}

\Verse{12}
{
	\hE{וַיֹּאמֶר יְהֹוָה אֶל מֹשֶׁה עֲלֵה אֵלַי הָהָרָה וֶהְיֵה שָׁם וְאֶתְּנָה לְךָ אֶת לֻחֹת הָאֶבֶן וְהַתּוֹרָה וְהַמִּצְוָה אֲשֶׁר כָּתַבְתִּי לְהוֹרֹתָם׃}
}
{
	\eE{
		And YHWH said to Moses, “Come up to me, to the mountain,
		and be there, and I’ll give you stone tablets and the
		instruction and the commandment that I’ve written to
		instruct them.”
	}
}
{
}

\Verse{13}
{
	\hE{וַיָּקם מֹשֶׁה וִיהוֹשֻׁעַ מְשָׁרְתוֹ וַיַּעַל מֹשֶׁה אֶל הַר הָאֱלֹהִים׃}
}
{
	\eE{
		And Moses and Joshua, his attendant, got up, and Moses
		went up to the Mountain of {\God}.
	}
}
{
}

\Verse{14}
{
	\hE{וְאֶל הַזְּקֵנִים אָמַר שְׁבוּ לָנוּ בָזֶה עַד אֲשֶׁר נָשׁוּב אֲלֵיכֶם וְהִנֵּה אַהֲרֹן וְחוּר עִמָּכֶם מִי בַעַל דְּבָרִים יִגַּשׁ אֲלֵהֶם׃}
}
{
	\eE{
		And he said to the elders, “Sit for us here until we’ll
		come back to you. And here, Aaron and Hur are with you.
		Let whoever has any matters go over to them.”
	}
}
{
}

\Verse{15}
{
	\hE{וַיַּעַל מֹשֶׁה אֶל הָהָר}
	\hP{וַיְכַס הֶעָנָן אֶת הָהָר׃}
}
{
	\eE{And Moses went up to the mountain.}
	\eP{And the cloud covered
		the mountain.}
}
{
}

\Verse{16}
{
	\hP{וַיִּשְׁכֹּן כְּבוֹד יְהֹוָה עַל הַר סִינַי וַיְכַסֵּהוּ הֶעָנָן שֵׁשֶׁת יָמִים וַיִּקְרָא אֶל מֹשֶׁה בַּיּוֹם הַשְּׁבִיעִי מִתּוֹךְ הֶעָנָן׃}
}
{
	\eP{
		And YHWH’s glory settled on Mount Sinai, and the cloud
		covered it six days. And He called to Moses on the seventh
		day from inside the cloud.
	}
}
{
}

\Verse{17}
{
	\hP{וּמַרְאֵה כְּבוֹד יְהֹוָה כְּאֵשׁ אֹכֶלֶת בְּרֹאשׁ הָהָר לְעֵינֵי בְּנֵי יִשְׂרָאֵל׃}
}
{
	\eP{
		And the appearance of YHWH’s glory was like a consuming
		fire in the mountaintop before the eyes of the children of
		Israel.
	}
}
{
}

\Verse{18}
{
	\hP{וַיָּבֹא מֹשֶׁה בְּתוֹךְ הֶעָנָן וַיַּעַל}
	\hE{אֶל הָהָר וַיְהִי מֹשֶׁה בָּהָר אַרְבָּעִים יוֹם וְאַרְבָּעִים לָיְלָה׃}
}
{
	\eP{And Moses came inside the cloud and went up}
	\eE{into the
		mountain. And Moses was in the mountain forty days and
		forty nights.}
}
{
}
